\documentclass[preview]{standalone}

\usepackage{amsmath}
\usepackage{amssymb}
\usepackage{stellar}
\usepackage{definitions}

\begin{document}

\id{psicologia-cognizione-processi}
\genpage

\section{Cognizione e processi mentali}

\begin{snippetdefinition}{cognizione-definition}{Cognizione}
    La \emph{cognizione} è un termine generale che si riferisce a tutte le forme
    di conoscenza. Lo studio della cognizione coincide quindi con lo studio
    della vita mentale.
\end{snippetdefinition}

\begin{snippet}{contenuti-processi-cognizione}
    La cognizione comprende contenuti e processi. I \emph{contenuti} sono ciò
    che si sa, come concetti, fatti, regole e ricordi. I \emph{processi
    cognitivi} sono invece i modi in cui questi contenuti vengono elaborati per
    interpretare il mondo e trovare soluzioni ai problemi.
\end{snippet}

\includesnpt[width=62\%,src=/snippet/static-psychology/cognitive-psychology-domains.jpg]{centered-img}

\begin{snippetdefinition}{psicologia-cognitiva-definition}{Psicologia cognitiva}
    La \emph{psicologia cognitiva} è il settore della psicologia che studia la
    cognizione e i processi mentali.
\end{snippetdefinition}

\begin{snippetdefinition}{scienze-cognitive-definition}{Scienze cognitive}
    Le \emph{scienze cognitive} raccolgono e integrano il sapere proveniente da
    discipline diverse che studiano la cognizione e la mente.
\end{snippetdefinition}

\includesnpt[width=48\%,src=/snippet/static-psychology/cognitive-science-interdisciplinarity.jpg]{centered-img}

\begin{snippet}{problema-studio-cognizione}
    Studiare la cognizione è difficile perché i processi cognitivi non sono
    direttamente osservabili. Si possono osservare l'input, come un messaggio
    ricevuto, e l'output, come una risposta o un'azione; il problema scientifico
    consiste nel ricostruire i processi e le rappresentazioni che collegano
    l'uno all'altro.
\end{snippet}

\begin{snippetexample}{donders-processi-mentali-example}{Compiti di Donders}
    Donders studiò la velocità dei processi mentali confrontando compiti che
    richiedevano passaggi cognitivi diversi. In un compito, il partecipante
    doveva stabilire se una lettera fosse maiuscola o minuscola e tracciare una
    \texttt{C} se era maiuscola. In un secondo compito, doveva anche stabilire
    se la lettera maiuscola fosse una vocale o una consonante e tracciare
    rispettivamente \texttt{V} o \texttt{C}. Il maggiore tempo richiesto dal
    secondo compito veniva interpretato come indice dei passaggi mentali
    supplementari.
\end{snippetexample}

\includesnpt[width=78\%,src=/snippet/static-psychology/donders-mental-processes.jpg]{centered-img}

\begin{snippetdefinition}{processi-seriali-definition}{Processi seriali}
    I \emph{processi seriali} sono processi mentali che si svolgono uno dopo
    l'altro, secondo una sequenza ordinata.
\end{snippetdefinition}

\begin{snippetdefinition}{processi-paralleli-definition}{Processi paralleli}
    I \emph{processi paralleli} sono processi mentali che si sovrappongono
    nell'unità di tempo e avvengono simultaneamente.
\end{snippetdefinition}

\begin{snippetexample}{processi-seriali-paralleli-example}{Processi seriali e paralleli}
    Al ristorante, scegliere che cosa ordinare valutando una portata alla volta
    è un esempio di processo seriale. Quando invece il cameriere prende
    l'ordine, i processi che permettono di comprendere la domanda possono
    svolgersi insieme a quelli che preparano la risposta verbale: in questo caso
    i processi sono paralleli.
\end{snippetexample}

\includesnpt[width=68\%,src=/snippet/static-psychology/serial-parallel-processes.jpg]{centered-img}

\begin{snippetdefinition}{risorse-elaborazione-definition}{Risorse di elaborazione}
    Le \emph{risorse di elaborazione} sono risorse cognitive limitate che devono
    essere distribuite tra i diversi compiti mentali in atto. La loro
    distribuzione dipende dai processi attentivi.
\end{snippetdefinition}

\begin{snippetdefinition}{processi-controllati-definition}{Processi controllati}
    I \emph{processi controllati} richiedono attenzione e molte risorse
    cognitive. Nella maggior parte dei casi non è possibile svolgere più
    processi controllati contemporaneamente.
\end{snippetdefinition}

\begin{snippetdefinition}{processi-automatici-definition}{Processi automatici}
    I \emph{processi automatici} richiedono poca attenzione e possono spesso
    essere eseguiti insieme ad altri processi automatici senza produrre
    interferenze rilevanti.
\end{snippetdefinition}

\begin{snippetexample}{camminare-conversare-example}{Camminare e conversare}
    In condizioni normali è semplice camminare e conversare, perché navigazione
    spaziale e linguaggio possono procedere in parallelo. Se però il marciapiede
    è pieno di ostacoli, scegliere il percorso richiede più risorse e può
    interrompere momentaneamente la conversazione.
\end{snippetexample}

\begin{snippetexample}{stroop-numerico-example}{Interferenza numerica}
    Se viene chiesto di individuare quale tra due numeri sia fisicamente più
    grande, il compito diventa più difficile quando la grandezza fisica è
    incongruente con il valore numerico. Per esempio, se il \texttt{5} è scritto
    più grande del \texttt{7}, il riconoscimento automatico del valore numerico
    interferisce con il compito percettivo.
\end{snippetexample}

\includesnpt[width=52\%,src=/snippet/static-psychology/numeric-stroop-interference.jpg]{centered-img}

\end{document}
